\section{Conclusion}
\label{Sec:Conclusion}


We have presented an interpolation-based approach to transition 
invariant generation for proving and disproving termination of 
infinite-state systems. 
The approach extends safety-based nontermination analysis with a termination proving capability, 
unifying the two within a single framework. 
Our key insight is to exploit terminating traces produced by SNA
as a source of structural information, using Craig interpolation 
to construct well-founded transition invariants guided by the 
termination conditions. 
The interaction between termination and nontermination reasoning is central to our technique: 
transition invariants focus the nontermination analysis on not covered states,
while SNA provides terminating traces that drive invariant generation.
% Grounding the approach in safety verification provides a robust foundation, enabling 
% efficient checking of both termination and nontermination via 
% existing safety solvers. 
We proved the soundness of the algorithm and conducted an extensive
experimental evaluation. 
The results demonstrate the benefit of combining termination and nontermination analyses,
and show that our approach is particularly strong at proving nontermination and complementary to state-of-the-art tools,
solving unique instances, including two that no participant of the Termination Competition~2025 solved.



% \textbf{Acknowledgements.} This work was supported in part by SNSF grant No. 200021-236601.
% We thank Marek Jankola and Dirk Beyer for fruitful discussions 
% on termination analysis. 
% We also want to thank Konstantin's girlfriend - Lea Petrivalska, for moral support and discussions during paper's submission.